\chapter{教育と人材：分流、選抜、定着}

\section{離れる前の振り返り}

今日はシンガポールでの最終日で、明日はアメリカに戻る。

この残り少ない時間を使って、滞在中に観察したシンガポールの教育システムと人材メカニズムについての考えをまとめておきたい。まだ十分に体系化されていないかもしれないが、生の記録として、後で読み返したときに新たな理解が生まれるかもしれない。

\section{分流：制度化された早期選抜}

シンガポールの教育システムは、幼稚園、小学校、中学校、ポリテクニック（理工学院）、ジュニアカレッジ（高校）、大学という段階に大まかに分けられる。多くの国との違いは、シンガポールが比較的早い段階から生徒の分流を始めることだ。通常は中学校修了後に、成績と能力に応じて異なる進路に導かれる。

一部の生徒は「ポリテクニック」と呼ばれる学校に進む。これらは応用型・技術型教育を中心としており、多くの国の「専門学校」とは本質的に異なる――シンガポールのポリテクニックはかなり近代的で、カリキュラムは産業ニーズと密接に結びついており、バイオエンジニアリング、金融管理、クラウドコンピューティングなど先端分野をカバーしている。

私は南洋理工学院（Nanyang Polytechnic）について学ぶ機会があった。初めてその名を聞いた多くの人は、南洋大学（Nanyang University）と何らかの関係があると思いがちだが、実際には主に中学校卒業者を対象に招生し、職業技術学校に近い位置付けだ。キャンパスは特に広大ではないが、至るところに高技術に関連する言葉や設備が見られ、かなり先端的な印象を受ける。ここの卒業生のほとんどはそのまま就職し、各産業の技術者や現場エンジニアとなる。

\section{頂点：NUSとNTU}

大学レベルでは、シンガポールを代表する二校が国立シンガポール大学（NUS）と南洋理工大学（NTU）だ。両校とも国際大学ランキングで上位に名を連ね、卒業生の就職見通しは総じて良好だ。外国人学生にとっては、この二校の学位を持つことが、シンガポールの永住権（PR）申請において明確な優位性をもたらすことが多い。

南洋理工大学のキャンパスは丘の上に位置し、海から近く、自然環境が心地よい。歴史的には、シンガポールの高等教育システムの基礎はイギリス植民地時代にさかのぼる。NUSは当初から英語での教育を主としていた一方、NTUの前身は中国語を教学言語とする大学であり、数十年の政策変化を経て、最終的に英語教育体系へと転換した。

この転換自体が、シンガポールの言語政策における全体的な方向性を反映している――英語こそが、主流エリート層への入場券なのだ。

\section{言語：エリート層への敷居}

シンガポールでは、標準的な英語（英米規範に近いもの）がエリート層に入るための重要なツールだ。しかし日常生活では、ほとんどの人が規範的な英語を話すのではなく、ローカル化された混合言語――シングリッシュ（Singlish、シンガポール英語）を話す。

シングリッシュは英語、華語（特に福建語や広東語などの方言）、マレー語、タミル語の要素を融合させ、独特の文法構造と発音の特徴を形成している。政府は「正しい英語を話そう」運動を何度も推進してシングリッシュの使用を抑制しようとしてきたが、この言語はすでに民間に深く根付き、シンガポールの文化的アイデンティティの一部となっている。

ここでは、どんな英語を話すかが、他者があなたをどう定義するかをしばしば決める。

\section{人材政策：精密に設計された選抜メカニズム}

人材の誘致と定着に関するシンガポールの政策設計も、同じく一貫した精緻さを示している。

一方で、優秀な外国人学生――特に博士課程の研究生――に対して、シンガポールはかなり充実した奨学金支援を提供することが多い。大幅な学費減免や全額免除、そして生活費補助が付いている。このメカニズムは、低コストで高い可能性を持つ人材の種を引き寄せ、国内の学術・産業エコシステムに組み込むことを目的としている。

一方で、一般的な外国人学生の場合、シンガポールで学ぶには通常、地元学生よりもはるかに高い学費を払う必要がある。この差別化された価格設定は政府に相当な収入をもたらすと同時に、暗黙の参入障壁を形成している――一定の経済力または学力を持つ外国人学生だけがこのシステムに入ることができる。

移民政策の面では、シンガポールの地元大学、特に一流大学の学位を持つ申請者は、永住権や市民権の申請において明確な優位性を享受することが多い。こうして明確な経路が形成される――教育システムを通じて人材を選抜し、奨学金を通じて人材を引き寄せ、移民政策を通じて人材を定着させる。

\section{まとめ：自律した人材生産システム}

総じて、シンガポールは単一の手段に頼るのではなく、教育分流、学費制度、ビザ・移民政策を有機的に組み合わせ、比較的完全な人材育成・定着システムを構築している。

このシステムは機能的に非常に自律している――国が必要とする各種人材を育成し、高度な研究と産業に必要な国際的力を取り込み、ある程度まで人口の構造と質を管理している。

この設計は真剣に考える価値がある。
